\documentclass[11pt]{article}
\usepackage[a4paper,margin=28mm]{geometry}
\usepackage{amsmath,amssymb,amsthm,booktabs,array}
\usepackage[T1]{fontenc}
\usepackage[colorlinks=true,linkcolor=blue,citecolor=blue,urlcolor=blue]{hyperref}
\newtheorem{theorem}{Theorem}
\newtheorem{lemma}[theorem]{Lemma}
\newtheorem{corollary}[theorem]{Corollary}
\DeclareMathOperator{\Hom}{Hom}
\newcommand{\TAS}{\mathrm{TAS}}
\newcommand{\ph}{\varphi}
\title{Convergence of Chen and Lin's direction-flip threshold}
\author{Haobo Ma\\\small ChronoAI Pte Ltd, Singapore
\and Wenlin Zhang\\\small National University of Singapore\\\small ChronoAI Pte Ltd, Singapore}
\date{}
\begin{document}
\maketitle
\begin{abstract}
We answer Chen and Lin's Question 6.1 in Section 6 of
\emph{Oriented Paths with Few Direction Flips Are Tournament Anti-Sidorenko}
(arXiv:2609.14655v1) affirmatively: their strict flip threshold divided by
path length converges, by reflection and amplification of any finite
counterexample followed by bounded padding in the same tournament.
The limit is the infimum of the normalized flip counts, with one extra flip,
of all non-anti-Sidorenko paths. A finite witness gives an explicit eventual
upper bound; combined with Chen and Lin's lower bound, this yields a limit
between $1/1665$ and $9/19$. Its exact value remains undetermined.
\end{abstract}

\section{The threshold and its limit}
A tournament is a finite loopless directed graph with exactly one arc between
each pair of distinct vertices. All tournaments below have at least one vertex.
For $k\geq1$ and $w=(w_1,\ldots,w_k)\in\{+,-\}^k$, let $P(w)$ have vertices
$0,\ldots,k$, with $i-1\to i$ when $w_i=+$ and $i\to i-1$ when $w_i=-$.
Write
\[
 \ph(w)=\#\{i:1\leq i<k,\ w_i\ne w_{i+1}\}.
\]
For an $N$-vertex tournament $T$, set
$t(w,T)=|\Hom(P(w),T)|/N^{k+1}$, where a homomorphism preserves arcs
and need not be injective. The word $w$ is tournament anti-Sidorenko,
abbreviated TAS, if $t(w,T)\leq2^{-k}$ for every such $T$.
Following Chen and Lin~\cite{CL}, define the \emph{strict} threshold
\begin{equation}\label{eq:threshold}
 f(k)=\max\{m\in\{0,\ldots,k\}:\ \ph(w)<m\Longrightarrow w\text{ is TAS
 for all }w\in\{+,-\}^k\}.
\end{equation}
Thus $f(k)=k$ if every length-$k$ word is TAS; otherwise $f(k)$ is the
least flip count of a non-TAS word. Indeed a word attaining that least count
excludes every larger $m$, and all words with smaller count are TAS.

Chen and Lin prove that every length-$k$ word with $r$ flips is TAS when
$k\geq1665r+1454$~\cite[Theorem 1.3]{CL}. Their Corollary 1.4 implies
$\liminf f(k)/k\geq1/1665$; Question 6.1 asks whether the limit exists.
For one flip, He, Mani, Nie, Tung and Wei~\cite[Proposition 2.4]{HMNTW} verified lengths
$3$ through $7$ by polynomial expansions, and Chen, Clemen and Noel~\cite[Theorem 1.4]{CCN}
proved the TAS property at every length at least $3$, using algebraic and
entropy methods. These are uniform results about prescribed low-flip words.
Chen and Lin's proof controls all words with few flips. Here the complementary
approach starts with one violating word and propagates it to \emph{every}
sufficiently large length. Reflection makes its transfer positive
semidefinite, allowing arbitrary repetition; a choice between the two next
edge directions handles the remainder.

The amplification inequality is the square-matrix case of Banks, Harcharras,
Neuwirth and Ricard~\cite[Theorem 1.1]{BHNR}; we include a spectral proof.
Endpoint-conditioned reflection in a fixed tournament also appears in
Chen, Clemen and Noel~\cite[Lemma 3.1]{CCN}.
Luck's analysis of the periodic rise--fall pattern $++--$~\cite[Section 8.2]{Luck}
already implies $\limsup f(k)/k\leq1/2$: its pattern probabilities have
$k$th root tending to $1/1.875104\ldots>1/2$, and finite transitive
tournaments approximate these probabilities. The contribution here is the
propagation of an arbitrary finite obstruction to every sufficiently large
length, yielding convergence and a finite-obstruction formula. Applying that
formula to earlier permutation counts gives the sharper bound $9/19$ below.

\begin{theorem}\label{thm:main}
Let $\mathcal D=\{1\}\cup\{(\ph(w)+1)/|w|:\ |w|\geq1,\ w\text{ is not TAS}\}$
and $L=\inf\mathcal D$. Then
\[
 \lim_{n\to\infty}\frac{f(n)}n=L,
 \qquad \frac1{1665}\leq L\leq\frac9{19},
 \qquad \frac{f(n)}n\geq L-\frac1n\quad(n\geq1).
\]
More explicitly, suppose $w$ has length $k\geq1$, $r=\ph(w)$, and
$q=2^kt(w,T)>1$ in an $N$-vertex tournament. Then $N\geq2$. Put
\begin{equation}\label{eq:cutoff}
 c=\frac{N-1}{N},\qquad
 M=1+\left\lfloor\frac{k\log(1/c)}{\log q}\right\rfloor.
\end{equation}
For every $n\geq2kM$ there is a non-TAS word of length $n$ in this same $T$, and
\begin{equation}\label{eq:upper}
 \frac{f(n)}n\leq\frac{r+1}{k}+\frac{2k}{n}.
\end{equation}
\end{theorem}

\section{Finite-host amplification}
Fix $T$ on $[N]=\{1,\ldots,N\}$, and let $A_{ij}=2/N$ if $i\to j$,
and $A_{ij}=0$ otherwise. Write $A_+=A$, $A_-=A^{\mathsf T}$,
$B_w=A_{w_1}\cdots A_{w_k}$, and
$\mu(B)=\mathbf1^{\mathsf T}B\mathbf1/N$; $\mathbf1$ is the all-one column
vector. The empty product is $I$. Expanding matrix multiplication gives
\begin{equation}\label{eq:mass}
 \mu(B_w)=2^kt(w,T).
\end{equation}
Each index sequence is an arc-preserving map, with the factor $(2/N)^k/N$;
all vertex collisions permitted by the arcs are retained.
For $\bar +=-$ and $\bar -=+$, let
$w^*=(\bar w_k,\ldots,\bar w_1)$. Then $B_{w^*}=B_w^{\mathsf T}$.

\begin{lemma}\label{lem:amplification}
For a real $N\times N$ matrix $B$ with $\mu(B)>1$ and an integer $m\geq0$,
\[
 \mu((BB^{\mathsf T})^m)\geq \mu(B)^{2m}.
\]
\end{lemma}
\begin{proof}
Put $e=\mathbf1/\sqrt N$ and $C=BB^{\mathsf T}$. Cauchy--Schwarz gives
$e^{\mathsf T}Ce=\|B^{\mathsf T}e\|^2\geq(e^{\mathsf T}Be)^2=\mu(B)^2$.
By the spectral theorem, $C$ has nonnegative eigenvalues $\lambda_j$ and an
orthonormal eigenbasis $u_j$. The weights $p_j=(e^{\mathsf T}u_j)^2$ sum to one.
For $m\geq1$, convexity of $x\mapsto x^m$ on $[0,\infty)$ yields
\[
 e^{\mathsf T}C^me=\sum_jp_j\lambda_j^m
 \geq\left(\sum_jp_j\lambda_j\right)^m
 \geq\mu(B)^{2m}.
\]
For $m=0$ both sides are one.
\end{proof}

\begin{lemma}\label{lem:padding}
For any word $u$ and integer $s\geq0$, some word $v$ of length $s$ satisfies
\[
 \mu(B_{uv})\geq c^s\mu(B_u),\qquad c=(N-1)/N.
\]
\end{lemma}
\begin{proof}
The tournament identity gives $(A+A^{\mathsf T})\mathbf1=2c\mathbf1$.
Consequently
\[
 \mu(B_uA)+\mu(B_uA^{\mathsf T})=2c\mu(B_u).
\]
At least one of the two terms is at least $c\mu(B_u)$.
Choose that next sign, and repeat $s$ times. Induction applies because $c\geq0$.
\end{proof}

\begin{proof}[Proof of Theorem~\ref{thm:main}]
Suppose $w$ is non-TAS and choose $T$ with $q=\mu(B_w)>1$.
The one-vertex tournament has no positive-length homomorphisms, so $N\geq2$
and $0<c<1$. Let $W=ww^*$, of length $d=2k$.
For $m\geq1$ its $m$-fold concatenation obeys
\begin{equation}\label{eq:repeated}
 \mu(B_{W^m})\geq q^{2m},\qquad
 \ph(W^m)=2m(r+1)-1.
\end{equation}
The mass inequality is Lemma~\ref{lem:amplification}.
For the flip identity, $w^*$ has $r$ flips, the middle join of $ww^*$ is
always a flip, and the join of two copies of $W$ is also always a flip:
the last sign of $w^*$ is opposite the first sign of $w$.
Thus $m(2r+1)+(m-1)=2m(r+1)-1$.

Write $n=dm+s$ with $0\leq s<d$. If $n\geq dM$, then $m\geq M$ and
the strict choice in~\eqref{eq:cutoff} gives $c^dq^{2m}>1$.
Lemma~\ref{lem:padding} supplies $v$ of length $s$ with
\[
 \mu(B_{W^mv})\geq c^sq^{2m}\geq c^dq^{2m}>1.
\]
So $W^mv$ is non-TAS. Appending $s$ edges creates at most $s$ flips,
including the first join when $s>0$. Equation~\eqref{eq:threshold} now gives
\[
 f(n)\leq2m(r+1)-1+s.
\]
Since $n=2km+s$, $r+1\leq k$, and $s<2k$, division proves~\eqref{eq:upper}.

The set $\mathcal D$ is nonempty and contained in $[0,1]$, so $0\leq L\leq1$.
If $f(n)<n$, a non-TAS word with exactly $f(n)$ flips exists, hence
$L\leq(f(n)+1)/n$. If $f(n)=n$, the same inequality follows from $L\leq1$.
This proves the stated lower bound for each $n$, and
$\liminf f(n)/n\geq L$.
For every non-TAS $w$,~\eqref{eq:upper} gives
$\limsup f(n)/n\leq(\ph(w)+1)/|w|$; also $\limsup f(n)/n\leq1$.
Taking the infimum over $\mathcal D$ gives the reverse inequality.
Chen and Lin's Corollary 1.4 gives $L\geq1/1665$.
The upper bound $9/19$ is proved next.
\end{proof}

\section{A finite witness and the remaining question}
\begin{corollary}\label{cor:half}
For the threshold in~\eqref{eq:threshold}, $L\leq9/19<1/2$.
\end{corollary}
\begin{proof}
Let $a(n)$ count permutations of $n$ labels whose successive comparisons
are the length-$(n-1)$ prefix of $++--++--\cdots$.
The earlier table A058258~\cite{OEIS} gives
\[
 a(19)=908138776681,\qquad a(20)=13194844482843.
\]
For $u=(++--)^4++$, each of its $a(19)$ permutations has twenty extensions
by a final rank. Exactly $a(20)$ of these have the next periodic sign, $-$.
Thus the number of permutations with signature $w=u+=(++--)^4+++$ is
\[
 D=20a(19)-a(20)=4967931050777,
 \qquad 2^{19}D>20!.
\]
This word has nineteen edges and eight flips. In the transitive tournament
$T_N$ on $[N]$, each twenty-element subset supports exactly $D$ injective
homomorphisms. Therefore
\[
 t(w,T_N)\geq D\binom{N}{20}/N^{20}\longrightarrow D/20!>2^{-19}
 \quad\text{as }N\to\infty.
\]
Some finite $T_N$ violates the TAS bound, so
Theorem~\ref{thm:main} gives $L\leq(8+1)/19$.
The accompanying \texttt{verify.py} checks the two tabulated counts by
last-rank recurrence, checks $D$ independently by inclusion--exclusion on
descents, and verifies $2^{19}D\binom{4096}{20}>4096^{20}$.
\end{proof}

For a smaller illustrative host, the word $z=++--++--++$ already violates
TAS in $T_{128}$. Its permutation count $45541$ appears in Stein's
1975 letter~\cite[p. 2]{Stein}; the following finite-host count includes
noninjective homomorphisms as well. Set $b_0(j)=1$ for $1\leq j\leq128$,
and for $1\leq i\leq10$ set
\[
 b_i(j)=\begin{cases}
 \sum_{h<j}b_{i-1}(h),&z_i=+,\\
 \sum_{h>j}b_{i-1}(h),&z_i=-.
 \end{cases}
\]
Here $b_i(j)$ counts maps of the first $i$ edges ending at $j$, so
\[
 H=\sum_jb_{10}(j)=158125066372222954176.
\]
The integer difference is
\[
 2^{10}H-128^{11}=10804340513327658237952>0.
\]
Thus $z$ is non-TAS. The accompanying \texttt{verify.py} evaluates the
recurrence and independently sums the corresponding matrix product in
integer arithmetic.

\begin{center}
\small
\begin{tabular}{@{}p{.30\linewidth}p{.62\linewidth}@{}}
\toprule
Source item in Chen--Lin~\cite{CL} & Mathematical coverage\\
\midrule
Question 6.1 & Answered affirmatively by Theorem~\ref{thm:main}, for every
positive length and the original strict maximum.\\[3pt]
Problem 6.2 & The exact numerical value remains undetermined;
$\liminf f(k)/k=L\in[1/1665,9/19]$. The infimum formula characterizes $L$
but does not evaluate it.\\[3pt]
Section 6, high-flip Sidorenko condition & The negative result is already
proved there by Chen and Lin; it concerns the opposite inequality and is
not a consequence of the convergence argument.\\[3pt]
Corollary 1.4 & Its lower bound is cited here; the limit replaces its
lower limit without changing its constant.\\
One-flip question cited in Section 1 & The conjecture of He et al. is already
proved by Chen, Clemen and Noel~\cite{CCN}; it is not re-proved here.\\
\bottomrule
\end{tabular}
\end{center}

\paragraph{Verification scope.}
The pinned Lean file
\href{https://github.com/the-omega-institute/trureturing/blob/81968a52b83b3d79d4db8e60c34e05d953675dc0/D5/S3/Combinatorics/ExtremalDensity/ChenLinFlipLimit.lean}{\texttt{ChenLinFlipLimit.lean}}
contains \nolinkurl{source_fullConvergence}, whose proposition is convergence
of the original strict threshold on sign vectors over all finite nonempty
tournaments. The declarations \nolinkurl{sourceF_eq_f} and
\nolinkurl{mass_transfer_eq_density} identify the encodings and the normalized
homomorphism count. Its \nolinkurl{reflected_power_amplification},
\nolinkurl{tournament_padding}, and \nolinkurl{defect_propagation} supply the
finite-host mechanism. The explicit logarithmic cutoff, the exact flip identity
in~\eqref{eq:repeated}, and Corollary~\ref{cor:half} have the written proofs
and integer checks above. The finite checks supplement, rather than replace,
the all-length proof.

\begin{thebibliography}{7}
\bibitem{BHNR}
W. Banks, A. Harcharras, S. Neuwirth and \'E. Ricard,
\emph{Matrix inequalities with applications to the theory of iterated kernels},
Linear Algebra Appl. \textbf{362} (2003), 275--286.
\href{https://doi.org/10.1016/S0024-3795(02)00517-7}{doi:10.1016/S0024-3795(02)00517-7}.
\bibitem{CL}
H. Chen and Y. Lin,
\emph{Oriented Paths with Few Direction Flips Are Tournament Anti-Sidorenko},
arXiv:2609.14655v1 (2026), Theorem 1.3, Corollary 1.4, and Section 6,
Question 6.1 and Problem 6.2.
\href{https://arxiv.org/abs/2609.14655v1}{arXiv:2609.14655v1}.
\bibitem{HMNTW}
X. He, N. Mani, J. Nie, N. Tung and F. Wei,
\emph{New Sidorenko-type inequalities in tournaments},
arXiv:2512.11222v1 (2025), Proposition 2.4, p.~6.
\href{https://arxiv.org/abs/2512.11222v1}{arXiv:2512.11222v1}.
\bibitem{CCN}
H. Chen, F. C. Clemen and J. A. Noel,
\emph{On Tournament Anti-Sidorenko Orientations of Trees},
arXiv:2605.14138v2 (2026), Theorem 1.4, p.~2.
\href{https://arxiv.org/abs/2605.14138v2}{arXiv:2605.14138v2}.
\bibitem{Luck}
J. M. Luck,
\emph{On the frequencies of patterns of rises and falls},
Physica A \textbf{407} (2014), 252--275;
\url{https://arxiv.org/abs/1309.7764v2}.
\bibitem{OEIS}
The OEIS Foundation Inc.,
\emph{A058258: The 2-Up sequence},
\url{https://oeis.org/A058258}, accessed 17 November 2026.
\bibitem{Stein}
P. R. Stein, letter to N. J. A. Sloane, 17 March 1975, p. 2,
in \emph{Correspondence, 1975}, A005981,
\url{https://oeis.org/A005981/a005981.pdf}.
\end{thebibliography}
\bigskip
\noindent\small Copyright \textcopyright\ Haobo Ma and Wenlin Zhang.
Licensed under Creative Commons Attribution 4.0 International (CC BY 4.0).
\end{document}
